\section{Methods}\label{sec:methods}

\subsection{Arms (v0.3 four-arm matrix)}

The v0.2 adversarial review identified two distinct confounds in the v0.2 four-arm design: (a) any \texttt{haiku\_cascade} win could reflect the extra compute purchased by the cascade's $K$ candidates and second pass, not the architecture; and (b) it could reflect any 2-pass revise protocol, not the specific content of the \iast{vimar\'sa} brief. v0.3 redesigns the matrix around two new control arms that target these confounds directly:

\begin{enumerate}[leftmargin=1.4em]
\item \textbf{\texttt{haiku\_bare} (no PCE).} Anthropic Claude Haiku invoked through the v0.3 \emph{clean Haiku CLI substrate} (\S\ref{sec:engine}): subprocess flag isolation, scrubbed \texttt{HOME}, per-item \texttt{IntegrityProbe}, single sampler call. Architecture-vs-nothing primary control.
\item \textbf{\texttt{haiku\_cascade} (with PCE).} Same Haiku endpoint through the v0.3 cascade with \texttt{commit\_policy="event\_gated"}, $K{=}3$, \texttt{max\_tokens}$=200$, the active-inference uplift (aspect-conditioned BMR, Hopfield warm-start, plumbed \iast{cit}-temperature, per-item free-energy budget), and an always-shadow revision that lets H8.v3 be measurable on every item.
\item \textbf{\texttt{haiku\_bare\_2K\_scorer} (no PCE, +K compute).} Single pass with \iast{icch\=a} fanning out $K' = 2K$ candidates and the same \iast{j\~n\=ana} BMR scorer the cascade uses (\texttt{commit\_policy="always\_draft"}). Spends roughly the same Haiku-token budget as the cascade but denies it the second pass and the \iast{vimar\'sa} brief. \emph{Extra-compute control} (H6.v3).
\item \textbf{\texttt{haiku\_generic\_revise\_2pass} (with 2-pass, no \iast{vimar\'sa} brief).} Two passes with \texttt{commit\_policy="always\_revise"} but the \iast{vimar\'sa}-generated brief is replaced by a fixed generic creative-revise prompt (\texttt{GENERIC\_REVISE\_BRIEF}). \emph{Revision-protocol control} (H7.v3).
\end{enumerate}
The \emph{primary v0.3 contrast} is \texttt{haiku\_cascade} vs.\ \texttt{haiku\_bare} (H1.v3--H4.v3 per domain). The \emph{headline fairness contrast} is \texttt{haiku\_cascade} vs.\ \texttt{haiku\_bare\_2K\_scorer} (H6.v3); the \emph{protocol fairness contrast} is \texttt{haiku\_cascade} vs.\ \texttt{haiku\_generic\_revise\_2pass} (H7.v3); and the \emph{within-cascade revision-causality contrast} is paired \texttt{score(revision)} -- \texttt{score(draft)} on items where the event-gated commit chose revision (H8.v3). The legacy \texttt{local\_*} arms are kept callable for backward-compatibility audits but are out of scope for v0.3 per the locked plan. All arms route through the same \texttt{GeneratorProtocol} so the only thing that varies between arms is the protocol, the compute budget, or the revision-brief content.

\subsection{Pre-registered hypotheses (v0.3)}

Eight hypotheses are pre-registered in \texttt{docs/SPEC\_v0.3.md} \S3 \emph{before} the v0.3 pilot is run:

\begin{itemize}[leftmargin=1.4em]
\item \textbf{H1.v3.} \texttt{haiku\_cascade} achieves higher CreativityPrism aggregate than \texttt{haiku\_bare} on the AUT slice.
\item \textbf{H2.v3.} \texttt{haiku\_cascade} achieves higher aspect-multiplicity score than \texttt{haiku\_bare} on the Wittgenstein aspect-shift poetry-interpretation probe.
\item \textbf{H3.v3.} \texttt{haiku\_cascade} achieves higher POEMetric advanced-creative-abilities composite than \texttt{haiku\_bare} on the poetry-generation slice.
\item \textbf{H4.v3.} \texttt{haiku\_cascade} achieves higher creative-novelty score than \texttt{haiku\_bare} on the BBH-style scientific-creativity probes.
\item \textbf{H5.v3 (redesigned).} The pooled effect across H1.v3--H4.v3, computed as a random-effects DerSimonian--Laird meta-analysis on the per-domain Hedges' $g$ values, is positive (with reported $\tau^2$ heterogeneity).
\item \textbf{H6.v3 (headline fairness).} \texttt{haiku\_cascade} achieves higher per-domain composite than \texttt{haiku\_bare\_2K\_scorer}. Architecture-vs-more-compute test.
\item \textbf{H7.v3 (protocol fairness).} \texttt{haiku\_cascade} achieves higher per-domain composite than \texttt{haiku\_generic\_revise\_2pass}. Brief-content vs.\ pass-existence test.
\item \textbf{H8.v3 (revision-causality).} Within-cascade paired \texttt{score(surface\_revision)} $-$ \texttt{score(surface\_draft)} on items where the event-gated commit chose revision is positive.
\end{itemize}

\subsection{Statistical protocol}

For each hypothesis we report: paired mean $\Delta$ (raw and length-controlled, where the per-arm linear word-count effect is regressed out before pairing), Hedges' $g$ with small-sample correction (raw and length-controlled), paired permutation $p$ (one-sided directional, exact for $n\le 18$, Monte-Carlo with $50000$ permutations otherwise), Wilcoxon signed-rank $p$ (one-sided), 95\% BCa bootstrap CI of the paired mean ($10000$ resamples), Holm--Bonferroni adjusted $p$ across $\{H_1, H_2, H_3, H_4\}$, and a-priori (assumed $g{=}0.5$) and retrospective (observed $g$) statistical power. A hypothesis is \emph{supported} iff its Holm-adjusted $p<0.05$ \emph{and} its BCa CI is strictly positive. H5.v3 is reported as a single random-effects pool over the four primary contrasts; H6.v3 and H7.v3 are reported per-domain (mirroring H1.v3--H4.v3) plus a meta-aggregated random-effects pool. H8.v3 is a within-arm paired test on items where the cascade's own commit policy chose revision. Every JSON leaf passes through \texttt{\_clean\_json}: non-finite floats (\texttt{NaN}, $\pm$\texttt{inf}) become \texttt{null} and the writer uses \texttt{allow\_nan=False}, so \texttt{benchmarks/results\_v0.3/stats.json} is RFC-compliant strict JSON (a v0.2 review finding).

\subsection{Reproducibility}

All seeds are recorded; bootstrap and permutation RNGs are seeded by \texttt{--seed} (default 4242). The benchmark driver (\texttt{benchmarks/driver.py}) writes a checkpoint after every call so a re-run resumes exactly from the failure point and runs the per-item \texttt{IntegrityProbe} (cached on \texttt{(env\_hash, flags\_hash)}; halts on leakage unless \texttt{--allow-leakage} is set). The v0.3 audit log lives under \texttt{audit/v0.3/} (with \texttt{audit/v0.3/integrity\_probes.jsonl}, \texttt{audit/v0.3/cost\_snapshot.json}, and \texttt{audit/v0.3/pilot\_run.log}); the legacy v0.2 audit lives under \texttt{audit/} (preserved unchanged). Every numerical claim in this paper is generated from \texttt{benchmarks/results\_v0.3/stats.json} via \texttt{benchmarks/autoreport.py} and is therefore traceable end-to-end.

\subsection{TRIZ-resolved contradictions (v0.3)}

The v0.2 adversarial review (\texttt{docs/reviews/2026-04-29-adversarial-v0.2-review.md}) identified five operator/protocol-level contradictions whose unresolved interactions explained the v0.2 confounds. v0.3 resolves each via TRIZ contradiction analysis (\texttt{docs/triz/v0.3/}) and a corresponding ADR (\texttt{docs/adr/v0.3/}):

\begin{enumerate}[leftmargin=1.4em]
\item \textbf{C1 / ADR-001 (clean substrate vs.\ OAuth).} The Haiku CLI authenticates via the macOS keychain (or \texttt{\textasciitilde/.config/claude/} on Linux); we cannot strip \texttt{HOME} entirely without losing auth. Resolved by selectively symlinking the keychain into a scrubbed temporary \texttt{HOME}, denying access to all other Claude Code state, and verifying the inner subprocess sees no plugins/skills/MCP via \texttt{IntegrityProbe} + \texttt{LEAKAGE\_REGEX}.
\item \textbf{C2 / ADR-002 (event-gated commit vs.\ guaranteed measurability).} Event-gated commit only records H8.v3 on items where the event fires; we want H8.v3 measurable on every item. Resolved by always generating a \emph{shadow} revision regardless of commit policy, so the paired \texttt{score(revision)} -- \texttt{score(draft)} test is well-defined for every item.
\item \textbf{C3 / ADR-003 (BMR aspect-conditioned reductions).} The v0.2 \iast{j\~n\=ana} reduces over a constant prior, producing degenerate $\Delta F$. Resolved by adding \texttt{ReductionTarget="aspect\_conditioned"}: when aspects are supplied, BMR reduces over $|\text{aspects}|$ candidate models (soft-boost in-support candidates), producing a non-degenerate $\Delta F$ that responds to aspect priors.
\item \textbf{C4 / ADR-004 (Hopfield in cascade vs.\ purity).} The v0.2 Hopfield store was only touched at consolidation, never on the causal path. Resolved by adding \texttt{ActiveHopfieldStore} (distinct from the legacy v0.1 store) read by \iast{apohana} as a warm-start aspect prior and written by \iast{vimar\'sa.consolidate} after every committed surface.
\item \textbf{C5 / ADR-005 (free-energy budget vs.\ open-ended cascade).} v0.2 had no per-item budget cap on the cascade's token spend. Resolved by introducing \texttt{FreeEnergyBudget}: a per-item ledger (\texttt{earn\_jnana}, \texttt{earn\_aspect}, \texttt{earn\_tokens}, \texttt{should\_continue\_revision}) that gates whether the revision pass runs at all under depleted budget.
\end{enumerate}
Each ADR is end-to-end traceable: TRIZ card $\to$ ADR $\to$ operator code $\to$ unit test $\to$ pilot result.
